% PhysicsLibrary submission source
% Suggested title: Mechanics learning path
% Suggested entry type: Topic
% Suggested primary PACS classification: 45.05.+x (General theory of classical mechanics)
% Suggested additional classifications: 45.20.Dd, 45.20.Jj, 45.40.-f, 45.50.-j, 45.50.Pk, 47.10.+g
% Suggested keywords: classical mechanics, Newtonian mechanics, analytical mechanics,
% Lagrangian mechanics, Hamiltonian mechanics, GRE physics, rigid body dynamics,
% oscillations, central force, non-inertial reference frame, fluid dynamics
% License intent: CC BY-SA 4.0

\section*{Mechanics learning path}

A \emph{mechanics learning path} is a prerequisite-ordered sequence of topics that organizes
classical mechanics from introductory particle motion through the analytical formulations used
in advanced undergraduate and graduate study.  Its purpose is pedagogical rather than taxonomic:
a classification scheme groups related subjects, whereas a learning path shows which ideas are
normally best understood before later ideas are attempted.

Classical mechanics is particularly well suited to this organization because its major methods
form a connected hierarchy.  A student first learns how to describe motion, then how interactions
change motion, and then how energy and momentum provide complementary descriptions of the same
dynamics.  Those ideas extend naturally to rotation, gravitation, oscillations, fluids, constrained
systems, rigid bodies, and finally the Lagrangian and Hamiltonian formulations.

The central progression is
\[
\begin{aligned}
\text{motion}
&\longrightarrow \text{forces}
\longrightarrow \text{energy and momentum}
\longrightarrow \text{rotation and systems} \\
&\longrightarrow \text{gravitation, oscillations, and fluids}
\longrightarrow \text{generalized coordinates} \\
&\longrightarrow \text{Lagrangian mechanics}
\longrightarrow \text{Hamiltonian mechanics}.
\end{aligned}
\]

\begin{figure}[h]
\centering
\setlength{\unitlength}{1mm}
\begin{picture}(118,112)
\put(14,100){\framebox(90,9){Mathematical foundations}}
\put(59,100){\vector(0,-1){7}}
\put(14,84){\framebox(90,9){Kinematics}}
\put(59,84){\vector(0,-1){7}}
\put(14,68){\framebox(90,9){Newtonian dynamics}}
\put(59,68){\vector(0,-1){7}}
\put(14,52){\framebox(90,9){Energy and momentum}}
\put(59,52){\vector(0,-1){7}}
\put(14,36){\framebox(90,9){Rotation, gravitation, oscillations, and fluids}}
\put(59,36){\vector(0,-1){7}}
\put(14,20){\framebox(90,9){Generalized coordinates and Lagrangian mechanics}}
\put(59,20){\vector(0,-1){7}}
\put(14,4){\framebox(90,9){Hamiltonian and graduate classical mechanics}}
\end{picture}
\caption{A prerequisite-ordered route from first-year mechanics to graduate classical mechanics.}
\end{figure}

This entry is intended to serve as a student-facing map of the mechanics portion of PhysicsLibrary.
The PACS classification remains useful for cataloging individual entries, but the order below is
chosen primarily according to conceptual dependence.

\section{Four natural entry points}

\subsection*{First-year calculus-based physics}

A beginning university student should start with mathematical foundations and kinematics, then
proceed through Newtonian dynamics, work and energy, momentum, rotation, gravitation, oscillations,
and elementary fluids.  This route supplies the conventional mechanics portion of a calculus-based
freshman physics course.

\subsection*{GRE Physics preparation}

A student preparing for the GRE Physics Subject Test should first verify mastery of introductory
mechanics and then review the topics explicitly included in the classical-mechanics portion of the
test.  ETS currently assigns about twenty percent of the test to classical mechanics.  The listed
scope includes kinematics, Newton's laws, work and energy, oscillatory motion, rotational motion,
systems of particles, central forces and celestial mechanics, three-dimensional particle dynamics,
Lagrangian and Hamiltonian methods, non-inertial frames, and elementary fluid dynamics.

\subsection*{Upper-undergraduate analytical mechanics}

A student beginning a junior- or senior-level mechanics course should already be comfortable with
Newton's laws, conservation laws, rotation, gravitation, ordinary differential equations, and
basic linear algebra.  The natural starting point is then generalized coordinates and constraints,
followed by Hamilton's principle, Lagrange's equations, central forces, small oscillations,
non-inertial frames, and rigid-body dynamics.

\subsection*{Graduate classical mechanics}

A graduate student should regard the earlier stages as review and continue through Hamiltonian
mechanics, Poisson brackets, canonical transformations, Hamilton--Jacobi theory, action-angle
variables, perturbation theory, advanced rigid-body motion, nonlinear dynamics, and continuum or
field extensions.

\section{Readiness diagnostic}

The following questions help locate an appropriate entry point.  They are not intended as a formal
examination.

\subsection*{Preparation for introductory mechanics}

A student should be able to answer yes to most of the following.
\begin{enumerate}
\item Can a vector be resolved into Cartesian components?
\item Can scalars and vectors be distinguished physically and mathematically?
\item Can a derivative be interpreted as an instantaneous rate of change?
\item Can an integral be interpreted as accumulated change or area under a curve?
\item Can elementary algebraic equations be solved for an unknown quantity?
\item Can basic trigonometric relations be used in mechanical geometry?
\end{enumerate}

\subsection*{Preparation for analytical mechanics}

Before beginning Lagrangian mechanics, a student should normally be able to:
\begin{enumerate}
\item draw and use a free-body diagram;
\item decide when Newton's second law, energy conservation, or momentum conservation is most efficient;
\item solve elementary collision problems;
\item compute or use a moment of inertia;
\item use torque and angular momentum as vectors;
\item solve the simple harmonic oscillator differential equation;
\item use partial derivatives and elementary multivariable functions;
\item carry out basic matrix operations.
\end{enumerate}

\subsection*{Preparation for graduate mechanics}

A student entering graduate classical mechanics should additionally be comfortable with generalized
coordinates, Lagrange's equations, multivariable Taylor expansions, linear algebra for coupled
differential equations, vector calculus, coordinate transformations, and first-order systems of
ordinary differential equations.

\section{Stage A: mathematical and kinematic foundations}

The first stage establishes the mathematical language used to describe motion before considering
what causes it.

\subsection*{Units and dimensional analysis}

Dimensional analysis provides an immediate check on equations and is often useful for estimating
the functional dependence of a result.  Mechanics students should be fluent with SI units, base
dimensions, scaling arguments, and order-of-magnitude reasoning.

\subsection*{Scalars and vectors}

Position, displacement, velocity, acceleration, force, momentum, torque, and angular momentum all
require careful use of vector algebra.  Essential tools include components, unit vectors, dot
products, and cross products.

\subsection*{Coordinate systems}

Cartesian coordinates should be learned first.  Cylindrical and spherical coordinates become
important when the geometry of a problem has rotational or radial symmetry.  Coordinate choice
becomes still more important later in generalized-coordinate mechanics.

\subsection*{Position, velocity, and acceleration}

The basic kinematic hierarchy is
\[
\mathbf{v}(t)=\frac{d\mathbf{r}}{dt},
\qquad
\mathbf{a}(t)=\frac{d\mathbf{v}}{dt}
             =\frac{d^2\mathbf{r}}{dt^2}.
\]
A student should be able to move in both directions through this hierarchy: differentiating a
trajectory to obtain velocity and acceleration, and integrating acceleration subject to initial
conditions to recover velocity and position.

\subsection*{Motion graphs and elementary trajectories}

Position--time, velocity--time, and acceleration--time graphs connect geometric intuition with
calculus.  Constant-acceleration equations should be understood as consequences of integration,
not merely memorized formulas.  Projectile motion, uniform circular motion, and relative motion
then extend the same ideas to two and three dimensions.

\paragraph{Checkpoint.}
Given a trajectory $\mathbf{r}(t)$, the student should be able to determine velocity and
acceleration, interpret their directions, and describe the resulting motion.

\section{Stage B: Newtonian dynamics}

Dynamics connects observed motion to physical interactions.  Newton's second law is
\[
\sum \mathbf{F}=\frac{d\mathbf{p}}{dt},
\]
which reduces to
\[
\sum \mathbf{F}=m\mathbf{a}
\]
for constant mass.

The central skill in elementary dynamics is converting a physical description into a force model.
This requires mastery of free-body diagrams and common interactions such as weight, normal force,
tension, static and kinetic friction, spring forces, and drag forces.  Connected bodies, pulleys,
inclined planes, and curved paths provide important examples because they require careful choices
of system boundary and coordinate axes.

Circular motion is important because constant speed does not imply zero acceleration.  For uniform
circular motion,
\[
a_r=\frac{v^2}{r}=\omega^2 r,
\]
directed toward the center of curvature.

\paragraph{Checkpoint.}
Given a mechanical configuration, the student should be able to isolate an appropriate body or
system, draw a free-body diagram, choose useful axes, and write the corresponding equations of
motion.

\section{Stage C: work, energy, and potentials}

Newton's laws give a local description through acceleration.  Work and energy often provide a more
direct global description between two states.  For a particle acted on by a force $\mathbf F$,
\[
W_{1\to2}=\int_1^2 \mathbf F\cdot d\mathbf r.
\]
The net work satisfies the work--energy theorem
\[
W_{\mathrm{net}}=\Delta T,
\qquad
T=\frac12 mv^2.
\]

For a conservative force a potential energy $U$ may be introduced so that
\[
\mathbf F=-\nabla U.
\]
In one dimension,
\[
F_x=-\frac{dU}{dx}.
\]
When only conservative forces do work, the mechanical energy $E=T+U$ is constant.

Potential-energy diagrams introduce turning points and stability.  An equilibrium position $q_0$
satisfies
\[
\left.\frac{dU}{dq}\right|_{q_0}=0,
\]
and a one-dimensional equilibrium is stable when
\[
\left.\frac{d^2U}{dq^2}\right|_{q_0}>0.
\]

\section{Stage D: momentum and systems of particles}

Linear momentum is $\mathbf p=m\mathbf v$ for a particle of constant mass.  Integrating Newton's
second law over time gives the impulse--momentum relation
\[
\mathbf J=\int_{t_1}^{t_2}\mathbf F\,dt=\Delta\mathbf p.
\]
For an isolated system, total linear momentum is conserved.  This is the natural tool for collisions
and explosions.

For a system of particles of total mass $M$, the center of mass is
\[
\mathbf R=\frac{1}{M}\sum_i m_i\mathbf r_i,
\]
and its motion obeys
\[
M\ddot{\mathbf R}=\sum \mathbf F_{\mathrm{ext}}.
\]
This demonstrates that an appropriate system boundary can eliminate internal forces from the
equations of motion.

\section{Stage E: elementary rotational mechanics}

Fixed-axis rotation parallels translational mechanics.  Torque is
\[
\boldsymbol\tau=\mathbf r\times\mathbf F,
\]
and for a rigid body rotating about a fixed principal axis,
\[
\sum\tau=I\alpha.
\]
The moment of inertia is
\[
I=\int r_\perp^2\,dm,
\]
and the rotational kinetic energy is
\[
T_{\mathrm{rot}}=\frac12 I\omega^2.
\]
Angular momentum satisfies
\[
\frac{d\mathbf L}{dt}=\boldsymbol\tau_{\mathrm{ext}}.
\]
Rolling without slipping combines translation and rotation through
\[
v_{\mathrm{cm}}=R\omega.
\]

\section{Stage F: gravitation and central forces}

Newtonian gravitation provides a major application in which energy and angular momentum work
together.  For point masses,
\[
\mathbf F=-\frac{GMm}{r^2}\hat{\mathbf r},
\qquad
U(r)=-\frac{GMm}{r}.
\]
Introductory study should include the gravitational field, escape speed, circular orbits, orbital
energy, and Kepler's laws.  At the analytical-mechanics level, the two-body problem is reduced to
one body with reduced mass
\[
\mu=\frac{m_1m_2}{m_1+m_2}.
\]
Conservation of angular momentum confines central-force motion to a plane and leads naturally to
effective potentials, orbit equations, and scattering.

\section{Stage G: oscillations}

The simple harmonic oscillator is a prototype for linear dynamics.  For
\[
m\ddot x+kx=0,
\]
the angular frequency is
\[
\omega_0=\sqrt{\frac{k}{m}},
\]
and the motion may be written
\[
x(t)=A\cos(\omega_0 t+\phi).
\]
The introductory sequence should include oscillator energy, the simple pendulum, damping, forced
motion, resonance, phase, bandwidth, and the quality factor.  The next step is coupled oscillation
and normal modes.

\section{Stage H: elementary fluids and elasticity}

Elementary fluids are part of a complete undergraduate mechanics foundation.  The sequence should
include density, pressure, hydrostatic equilibrium, Pascal's principle, Archimedes' principle, the
continuity equation, Bernoulli's equation, viscosity, and elementary pipe flow.

For a static fluid of density $\rho$ in a uniform gravitational field,
\[
\frac{dp}{dz}=-\rho g.
\]
For steady incompressible flow,
\[
A_1v_1=A_2v_2.
\]
Along a streamline for an ideal incompressible fluid under the usual Bernoulli assumptions,
\[
p+\frac12\rho v^2+\rho gz=\text{constant}.
\]
Elasticity provides a bridge from particle mechanics to continuum mechanics through stress, strain,
elastic moduli, and elastic energy.

\section{Stage I: generalized coordinates and constraints}

Analytical mechanics begins by asking what coordinates actually describe the independent motion of
a system.  A set of generalized coordinates $q_1,\ldots,q_n$ specifies the configuration using the
number of independent variables required by the constraints.

Important distinctions include holonomic and nonholonomic constraints and scleronomic and rheonomic
constraints.  Virtual displacements and d'Alembert's principle explain how ideal constraint forces
can often be eliminated from the equations.

\section{Stage J: variational and Lagrangian mechanics}

For many mechanical systems the Lagrangian is
\[
L(q_i,\dot q_i,t)=T-U.
\]
The action is
\[
S[q]=\int_{t_1}^{t_2}L(q_i,\dot q_i,t)\,dt.
\]
Hamilton's principle leads to the Euler--Lagrange equations
\[
\frac{d}{dt}\left(\frac{\partial L}{\partial \dot q_i}\right)
-\frac{\partial L}{\partial q_i}=0.
\]
A complete upper-undergraduate treatment should include generalized forces, Lagrange multipliers,
generalized momentum,
\[
p_i=\frac{\partial L}{\partial\dot q_i},
\]
cyclic coordinates, first integrals, and Noether's theorem.

\section{Stage K: non-inertial reference frames}

If a frame rotates with angular velocity $\boldsymbol\Omega$, derivatives measured in inertial and
rotating frames are related by the transport theorem
\[
\left(\frac{d\mathbf A}{dt}\right)_I
=
\left(\frac{d\mathbf A}{dt}\right)_R
+\boldsymbol\Omega\times\mathbf A.
\]
Applying the relation twice to position yields the centrifugal, Coriolis, and Euler terms that
appear in rotating-frame dynamics.  This topic is important for terrestrial motion, navigation,
geophysical dynamics, and rigid-body mechanics.

\section{Stage L: central-force mechanics and scattering}

Central-force motion combines conservation of energy and angular momentum with the effective
potential
\[
U_{\mathrm{eff}}(r)=U(r)+\frac{L^2}{2\mu r^2}.
\]
The progression should continue through the orbit equation, the Kepler problem, orbit stability,
impact parameter, scattering angle, differential cross section, and Rutherford scattering.

\section{Stage M: small oscillations and normal modes}

Near stable equilibrium, many systems can be linearized.  With generalized displacements
$\boldsymbol\eta$,
\[
\mathbf M\ddot{\boldsymbol\eta}+\mathbf K\boldsymbol\eta=0,
\]
which leads to the generalized eigenvalue problem
\[
\det(\mathbf K-\omega^2\mathbf M)=0.
\]
Normal modes provide a bridge from a single oscillator to molecular vibration, structural dynamics,
waves, and field modes.

\section{Stage N: three-dimensional rigid-body mechanics}

The rigid-body sequence should proceed through rotation representation, angular velocity, inertia
tensor, principal axes, angular momentum, Euler's equations, torque-free motion, precession, and
nutation.  In general,
\[
\mathbf L=\mathbf I\boldsymbol\omega,
\]
and $\mathbf L$ need not be parallel to $\boldsymbol\omega$.

\section{Stage O: Hamiltonian mechanics}

Canonical momenta are defined by
\[
p_i=\frac{\partial L}{\partial\dot q_i}.
\]
A Legendre transform gives the Hamiltonian
\[
H(q,p,t)=\sum_i p_i\dot q_i-L,
\]
and Hamilton's equations are
\[
\dot q_i=\frac{\partial H}{\partial p_i},
\qquad
\dot p_i=-\frac{\partial H}{\partial q_i}.
\]
The conceptual shift is from configuration-space second-order dynamics to phase-space first-order
dynamics.

\section{Stage P: Poisson and symplectic mechanics}

For functions $f(q,p,t)$ and $g(q,p,t)$, the Poisson bracket is
\[
\{f,g\}=\sum_i
\left(
\frac{\partial f}{\partial q_i}\frac{\partial g}{\partial p_i}
-
\frac{\partial f}{\partial p_i}\frac{\partial g}{\partial q_i}
\right).
\]
Poisson brackets package Hamiltonian time evolution, canonical structure, generators, and
conservation laws.  Liouville's theorem and symplectic geometry then reveal structural properties of
Hamiltonian flow.

\section{Stage Q: canonical transformations}

A canonical transformation changes phase-space variables while preserving Hamilton's form of the
equations.  Generating functions provide a systematic method for constructing these transformations.

\section{Stage R: Hamilton--Jacobi theory}

Hamilton--Jacobi theory seeks a canonical transformation that makes the transformed dynamics as
simple as possible.  Its central equation is
\[
H\left(q_i,\frac{\partial S}{\partial q_i},t\right)
+\frac{\partial S}{\partial t}=0.
\]
This theory also provides a conceptual bridge to geometrical optics and quantum mechanics.

\section{Stage S: action-angle variables and perturbation theory}

For integrable periodic systems, action-angle variables turn the motion into especially simple
canonical evolution.  They provide a natural language for adiabatic invariants and a convenient
starting point for perturbation theory in nearly integrable systems.

\section{Stage T: nonlinear dynamics and chaos}

Important ideas include fixed points, local stability, bifurcations, phase portraits, Poincare
sections, sensitivity to initial conditions, and deterministic chaos.  These topics show where the
simple picture of globally integrable motion breaks down and why qualitative and computational
methods become essential.

\section{Stage U: continuous systems and fields}

The variational structure of particle mechanics extends naturally to continuous systems.  Instead
of a finite-dimensional Lagrangian one introduces a Lagrangian density $\mathcal L$ and an action
\[
S=\int \mathcal L(\phi,\partial_\mu\phi,x)\,d^4x.
\]
The resulting field Euler--Lagrange equations provide a bridge from discrete mechanics to strings,
elastic media, fluids, and classical field theory.

\section{GRE mechanics checklist}

A student preparing specifically for the Physics Subject Test should verify competence in:
\begin{enumerate}
\item kinematics in one, two, and three dimensions;
\item Newton's laws and force modeling;
\item work, energy, and power;
\item linear momentum, impulse, and systems of particles;
\item oscillatory motion;
\item rotational motion about a fixed axis;
\item angular momentum and torque;
\item central forces, gravitation, and celestial mechanics;
\item three-dimensional particle dynamics;
\item Lagrangian and Hamiltonian formalism;
\item non-inertial reference frames;
\item elementary fluid dynamics.
\end{enumerate}

GRE preparation rewards both understanding and speed.  Students should therefore practice not only
full derivations but also dimensional reasoning, limiting cases, ratio arguments, conservation-law
shortcuts, and rapid recognition of standard equations.

\section{Recommended study method}

For each major topic, use a five-pass cycle.
\begin{enumerate}
\item Identify the physical system, degrees of freedom, interactions, constraints, and likely conservation laws.
\item Work through the governing derivation instead of memorizing an isolated formula.
\item Reproduce worked examples without looking at intermediate steps.
\item Solve conceptual, numerical, derivational, and computational exercises.
\item For GRE-relevant material, add short timed problems emphasizing recognition and efficient reasoning.
\end{enumerate}

Reading alone is not evidence of mastery.  The important skill is the ability to set up and solve a
new problem whose surface description differs from previously seen examples.

\section{Milestones}

\subsection*{Milestone 1: introductory particle mechanics}
The student can translate among words, graphs, and equations of motion; draw correct free-body
diagrams; solve Newton's-law problems; use energy and momentum appropriately; and analyze collisions
and center-of-mass motion.

\subsection*{Milestone 2: rotation, gravitation, oscillations, and fluids}
The student can solve fixed-axis rotational dynamics; use torque and angular momentum; analyze
rolling; solve elementary orbit problems; analyze harmonic oscillators, damping, and resonance; and
solve introductory hydrostatics, continuity, and Bernoulli problems.

\subsection*{Milestone 3: analytical mechanics}
The student can identify degrees of freedom and constraints, choose useful generalized coordinates,
derive Lagrange's equations, use generalized momenta and cyclic coordinates, and analyze central
forces, non-inertial frames, and small oscillations.

\subsection*{Milestone 4: rigid-body and Hamiltonian mechanics}
The student can use the inertia tensor and principal axes, formulate Euler's rigid-body equations,
construct a Hamiltonian, use Hamilton's equations, interpret phase space, and compute basic Poisson
brackets.

\subsection*{Milestone 5: graduate mechanics}
The student can work with canonical transformations and generating functions, solve representative
Hamilton--Jacobi problems, use action-angle variables, understand perturbative and nonlinear
departures from integrability, and connect discrete mechanics to continuum or field formulations.

\section{Source and licensing note}

This learning path is an original synthesis prepared for PhysicsLibrary.  Its ordering has been
cross-checked against openly licensed university mechanics resources and against the topic scope of
the GRE Physics Subject Test.

\begin{thebibliography}{9}

\bibitem{UCD9A}
T. Weideman,
\emph{UCD Physics 9A -- Classical Mechanics}, Physics LibreTexts.
The resource is distributed under a Creative Commons Attribution--ShareAlike license.
\PMlinkexternal{UCD Physics 9A}{https://phys.libretexts.org/Courses/University_of_California_Davis/UCD:_Classical_Mechanics}

\bibitem{WikibooksCM}
Wikibooks contributors,
\emph{Classical Mechanics}.
Wikibooks text is generally distributed under the Creative Commons Attribution--ShareAlike 4.0
license, subject to page-specific notices.
\PMlinkexternal{Wikibooks Classical Mechanics}{https://en.wikibooks.org/wiki/Classical_Mechanics}

\bibitem{OpenStax2016}
OpenStax,
\emph{University Physics, Volume 1}, archived 2016 release.
The archived edition states a Creative Commons Attribution 4.0 license except where otherwise noted.
\PMlinkexternal{Archived University Physics Volume 1}{https://pressbooks.bccampus.ca/universityphysicssandboxbook1/}

\bibitem{ETS}
Educational Testing Service,
\emph{GRE Physics Subject Test: Content and Structure}.
Used only to identify the scope and weighting of the classical-mechanics portion of the examination.
\PMlinkexternal{GRE Physics Subject Test content}{https://www.ets.org/gre/test-takers/subject-tests/about/content-structure.html}

\end{thebibliography}

Unless otherwise noted, this PhysicsLibrary entry is intended for release under the Creative Commons
Attribution--ShareAlike 4.0 International license.
